% !PRINTED 267
\weberpage{267}
%<<CONT
becomes real. But by (1) and (2) we obtain
\[
R\,S\,R^{-1} = \begin{pmatrix} \lambda\,\varrho\,\sigma - \mu\,\nu\,\sigma^{-1}, & -\lambda\,\mu\,(\sigma-\sigma^{-1}) \\ \varrho\,\nu\,(\sigma-\sigma^{-1}), & -\mu\,\nu\,\sigma+\varrho\,\lambda\,\sigma^{-1} \end{pmatrix},
\]
and this becomes real when $\lambda\mu$ and $\varrho\nu$ are purely imaginary, and $\lambda\varrho$ and
$-\mu\nu$ are conjugate imaginary. These conditions can be satisfied
in several ways: most simply, no doubt, and at the same time
in such a way that the determinant of $R = 1$, if one sets
\weq{4}{R = \begin{pmatrix} \frac{1}{2}, & -\frac{1}{2}\,\varepsilon^{-1} \\ \varepsilon, & 1 \end{pmatrix}}
If, on the other hand,
\[
A = \begin{pmatrix} a, & b \\ c, & d \end{pmatrix}
\]
is an arbitrary real substitution from $L_p$, then by (4),
since $\varepsilon^2 = N$, we obtain:
\[
R^{-1}\,A\,R =
\]
\weq{5}{\begin{pmatrix} \dfrac{a+d}{2}+b\,\varepsilon+\dfrac{c}{4}\,\varepsilon^{-1}, & -\dfrac{a-d}{2}\,\varepsilon^{-1}+b-\dfrac{c}{4}\,N^{-1} \\[1.2ex] -N\left(\dfrac{a-d}{2}\,\varepsilon^{-1}+b-\dfrac{c}{4}\,N^{-1}\right), & \dfrac{a+d}{2}-b\,\varepsilon-\dfrac{c}{4}\,\varepsilon^{-1} \end{pmatrix},}
for which we may also set
\weq{6}{A = \begin{pmatrix} \alpha, & \beta \\ -N\beta', & \alpha' \end{pmatrix} = \begin{pmatrix} \alpha, & \beta \\ -N\beta^p, & \alpha^p \end{pmatrix}, \quad \alpha\alpha'+N\beta\beta' = 1}
in which $\alpha,\alpha'$ and $\beta,\beta'$ are two conjugate imaginary pairs,
so that by §.\,65, (4) one may set $\alpha' = \alpha^p$, $\beta' = \beta^p$.

If, conversely, we take any substitution $A$ of the
form (6), then
\[
A = R\,A\,R^{-1} =
\]
\weq{7}{\begin{pmatrix} \dfrac{\alpha+\alpha'}{2}+\varepsilon\,\dfrac{\beta'-\beta}{2}, & \dfrac{\beta+\beta'}{4}+\varepsilon^{-1}\,\dfrac{\alpha-\alpha'}{4} \\[1.2ex] -N(\beta+\beta')+\varepsilon(\alpha-\alpha'), & \dfrac{\alpha+\alpha'}{2}-\varepsilon\,\dfrac{\beta'-\beta}{2} \end{pmatrix}}
results as a \gesperrt{real} substitution from $L_p$.

\gesperrt{From this it follows that, when $A$ runs through the whole group $L_p$,
the substitution $A$ determined by (5) and (6) runs through a group
isomorphic with $L_p$, which we denote by $\Lambda_p$}.
